\chapter{Operational Protocols \& Telemetry}
\label{chap:protocols}

Our distributed nodes run asynchronous telemetry listeners to monitor server status and coordinate transport actions in real time.

\section{Socket Configuration and Ports}
Communication channels strictly adhere to the internal network matrix documented in Table~\ref{tab:port_map}.

\begin{table}[htbp]
\centering
\caption{System Port Interface Matrix}
\label{tab:port_map}
\begin{tabular}{llll}
\toprule
\textbf{Port} & \textbf{Protocol} & \textbf{Traffic Type} & \textbf{Service Endpoint} \\ 
\midrule
25565 & TCP & Game Client Handshakes & Paper Minecraft Server \\
25575 & TCP & Administrative Console & RCON Control Bridge \\
5005  & TCP & Telemetry Events       & CART\_PASS Async Listener \\
\bottomrule
\end{tabular}
\end{table}

\section{Telemetry Event Processing}
When a physical minecart triggers an accelerated detector rail array on Chonk, the telemetry sender encapsulates metadata regarding the transit and forwards it to Stargate on Port 5005.

Let $E_{cart}$ be the telemetry event package. The serialized JSON envelope must strictly map to the following payload schema:

\begin{lstlisting}[style=bashstyle, caption={CART\_PASS Telemetry JSON Schema}]
{
  "event": "CART_PASS",
  "node": "CHONK-01",
  "data": {
    "rail_id": "WASHINGTON_TERM_04",
    "coordinates": {
      "x": -1042,
      "y": 64,
      "z": 1289
    },
    "timestamp": 1778346300
  }
}
\end{lstlisting}

\section{Debugging Pipelines}
If an AI agent confirms building generation but changes are not reflected in the live world, administrators must isolate the pipeline drop points:

\begin{equation}
    \text{T2BM Output} \xrightarrow{\text{Validated}} \text{test\_core.log} \xrightarrow{\text{Transited}} \text{hailort.log} \xrightarrow{\text{Injected}} \text{Chonk RCON}
\end{equation}

By sequentially running the unit test targets via \texttt{pytest}, drop points can be identified and cleared before pipeline timeouts occur.